\documentclass[11pt,a4paper]{article}
\usepackage[utf8]{inputenc}
\usepackage[T1]{fontenc}
\usepackage{amsmath,amssymb,amsthm}
\usepackage{bm}
\usepackage[margin=2.5cm]{geometry}
\usepackage{hyperref}

\newcommand{\uc}{\mathrm{uc}}
\newcommand{\Guc}{\Gamma_{\uc}}
\newcommand{\bdry}{\partial}
\newcommand{\inn}{\mathrm{in}}

\title{Gaussian fPEPS with larger unit cells:\\
construction, unified loss function and block-diagonalized Gaussian map}
\date{July 2026}

\begin{document}
\maketitle

\begin{abstract}
These notes document the extension of the GfPEPS translation scheme
from trivial $(1\times 1)$ unit cells to arbitrary
rectangular $(L_x \times L_y)$ unit cells, as implemented in
\texttt{GfPEPS.jl}. We review the construction of the unit-cell covariance
matrix $\Guc$ following Hackenbroich \emph{et al.}
[Phys.\ Rev.\ B \textbf{101}, 115134 (2020)], define the unified energy loss
function used for both BCS and Kitaev-type Hamiltonians, and derive the
\emph{block-diagonalization} of the Gaussian map obtained from a smart
ordering of the virtual modes: all momentum dependence is confined to the
virtual modes on bonds that wrap around the unit-cell boundary, so the
intra-cell modes can be eliminated once per loss evaluation by a
$k$-independent Schur complement. This reduces the per-momentum inversion
cost from $\mathcal{O}\!\big((8\Lambda L_x L_y)^3\big)$ to
$\mathcal{O}\!\big((4\Lambda (L_x{+}L_y))^3\big)$.
\end{abstract}

\section{Setup and conventions}

Each lattice site $i$ carries $N_f$ physical fermions and $4\Lambda$ virtual fermions
($\Lambda$ flavors on each of the four bonds $l,r,u,d$), described by the
local fiducial state $|Q_i\rangle$ with covariance matrix (CM)
\begin{equation}
  (\Gamma_Q)_{\mu\nu} = \tfrac{i}{2}\,
  \langle Q | [c_\mu, c_\nu] | Q \rangle ,
  \qquad
  \Gamma_Q =
  \begin{pmatrix} A & B \\ -B^{T} & D \end{pmatrix},
  \label{eq:GammaQ}
\end{equation}
where $A \in \mathbb{R}^{2N_f \times 2N_f}$ acts on the physical Majorana
modes, $D \in \mathbb{R}^{8\Lambda \times 8\Lambda}$ on the virtual ones,
and $B$ couples the two sectors. The virtual Majorana modes of each site are
ordered as
\begin{equation}
  \big(c^{(1)}_{l_1}, c^{(2)}_{l_1}, c^{(1)}_{r_1}, c^{(2)}_{r_1},
  \dots,
  c^{(1)}_{u_1}, c^{(2)}_{u_1}, c^{(1)}_{d_1}, c^{(2)}_{d_1}, \dots\big),
  \label{eq:siteordering}
\end{equation}
i.e.\ first the $2\Lambda$ pairs of left/right modes, then the $2\Lambda$
pairs of up/down modes. The physical state $|\Psi\rangle$ is obtained by
projecting all virtual bonds onto maximally entangled states
$|\omega_{ij}\rangle$; on the level of CMs this is the Gaussian map
\begin{equation}
  G^{(\Psi)}_{\bm k}
  = B \big( D + G^{(\omega)}_{\bm k} \big)^{-1} B^{T} + A ,
  \label{eq:gaussianmap}
\end{equation}
with $G^{(\omega)}_{\bm k}$ the Fourier-transformed CM of the maximally
entangled virtual bonds.

\section{Construction for an $L_x \times L_y$ unit cell}

For a translation-invariant state with an $L_x \times L_y$ unit cell we
assign one local fiducial state $|Q_s\rangle$ to every site
$s = i_x + (i_y - 1) L_x$ of the unit cell (sites related by the unit-cell
layout may share the same variational tensor). The fiducial state of the
whole unit cell is the product state of the per-site fiducial states, so its
CM is block diagonal,
\begin{equation}
  A_{\uc} = \bigoplus_{s=1}^{L_x L_y} A_s , \qquad
  B_{\uc} = \bigoplus_{s=1}^{L_x L_y} B_s , \qquad
  D_{\uc} = \bigoplus_{s=1}^{L_x L_y} D_s ,
  \label{eq:blockdiagABD}
\end{equation}
with the physical modes of all sites listed first (site-major ordering) and
the virtual modes of all sites following in the same site-major order with
the per-site ordering \eqref{eq:siteordering}. Each per-site block is
parametrized by an orthogonal matrix $X_s$ through the canonical form
$\Gamma_s = X_s^{T} \big(\bigoplus_{j} i\sigma_y\big) X_s$, exactly as in the
$1\times 1$ case, and the $X_s$ of all \emph{distinct} sites of the unit
cell are the variational parameters.

The virtual-bond CM factorizes into one term per nearest-neighbor bond. A
horizontal bond couples the $r$ modes of site $(i_x, i_y)$ to the $l$ modes
of site $(i_x{+}1, i_y)$; for $i_x < L_x$ both sites lie in the same unit
cell and the bond contributes the \emph{momentum-independent} block
\begin{equation}
  G^{(\omega)}\big[l_{(i_x+1,i_y)}, r_{(i_x,i_y)}\big] = -\sigma_x ,
  \qquad
  G^{(\omega)}\big[r_{(i_x,i_y)}, l_{(i_x+1,i_y)}\big] = +\sigma_x ,
\end{equation}
per flavor, whereas the wrapping bond $i_x = L_x$ connects to the
neighboring unit cell and acquires the phase of one unit-cell translation,
\begin{equation}
  G^{(\omega)}_{\bm k}\big[l_{(1,i_y)}, r_{(L_x,i_y)}\big]
  = -e^{i k_x L_x}\,\sigma_x ,
  \qquad
  G^{(\omega)}_{\bm k}\big[r_{(L_x,i_y)}, l_{(1,i_y)}\big]
  = +e^{-i k_x L_x}\,\sigma_x ,
  \label{eq:wrapbond}
\end{equation}
and analogously for vertical bonds with $e^{\pm i k_y L_y}$. For
$L_x = L_y = 1$ this reduces to Eq.~(C7) of \emph{Hackenbroich et al.}.
The allowed momenta live in the folded Brillouin zone
$k_\alpha \in [-\pi/L_\alpha, \pi/L_\alpha)$.

With \eqref{eq:blockdiagABD} and \eqref{eq:wrapbond}, the physical CM at
momentum $\bm k$ is again given by the Gaussian map
\eqref{eq:gaussianmap}, now of dimension $2 N_f L_x L_y$, obtained by
inverting the $8\Lambda L_x L_y$-dimensional matrix
$D_{\uc} + G^{(\omega)}_{\bm k}$ for every $\bm k$.

\section{Unified loss function}

The energy of any quadratic (BdG) Hamiltonian --- including the Kitaev
honeycomb model in the vortex-free sector, which maps onto a
$N_f = 1$ BCS-type Hamiltonian on the square lattice --- can be evaluated
directly from $G^{(\Psi)}_{\bm k}$. Writing the unit-cell Bloch Hamiltonian
in the Majorana basis as $h_{\bm k} = \tfrac{i}{2}\, \Omega H^{\mathrm{BdG}}_{\bm k} \Omega^\dagger$,
the energy per unit cell is
\begin{equation}
  E = \frac{1}{N_k} \sum_{\bm k} \Big[
      \tfrac12 \operatorname{tr} \xi_{\bm k}
      - \tfrac14 \operatorname{tr}\!\big( h_{\bm k}\, G^{(\Psi)\,T}_{\bm k} \big)
      + E_{\mathrm{shift}}(\bm k) \Big],
  \label{eq:loss}
\end{equation}
which is a single trace formula independent of the specific pairing or
hopping structure: BCS ($N_f = 2$) and Kitaev ($N_f = 1$, with
$E_{\mathrm{shift}} = -J_z L_x L_y$ accounting for the $\mathbb{Z}_2$ background gauge field). 
Hamiltonians only differ in the matrices $\xi_{\bm k}$, $\Delta_{\bm k}$ entering $H^{\mathrm{BdG}}_{\bm k}$. This
realizes the unified loss function: the same code path
\texttt{energy\_loss\_X} optimizes any \texttt{MomentumSpaceBdGHamiltonian}
on any rectangular unit cell.

\section{Block-diagonalization via smart mode ordering}
\label{sec:blockdiag}

The naive evaluation of \eqref{eq:gaussianmap} requires factorizing the
$d \times d$ matrix $D_{\uc} + G^{(\omega)}_{\bm k}$ with
$d = 8\Lambda L_x L_y$ \emph{for every momentum} $\bm k$, which dominated
the cost of previous attempts at larger unit cells. The key observation is
that the momentum dependence of $G^{(\omega)}_{\bm k}$ in
\eqref{eq:wrapbond} is carried \emph{only} by the wrapping bonds. We
therefore reorder the virtual Majorana modes into two groups:
\begin{align*}
  \bdry \ (\text{boundary}):&\quad
  l\text{-modes of column } 1,\;
  r\text{-modes of column } L_x,\;
  u\text{-modes of row } 1,\;
  d\text{-modes of row } L_y,\\
  \inn \ (\text{inner}):&\quad \text{all remaining virtual modes},
\end{align*}
with $d_\bdry = 4\Lambda(L_x + L_y)$ and $d_\inn = d - d_\bdry$. In this
ordering,
\begin{equation}
  D_{\uc} + G^{(\omega)}_{\bm k}
  =
  \begin{pmatrix}
    M_{\inn\inn} & D_{\inn\bdry} \\
    -D_{\inn\bdry}^{T} & D_{\bdry\bdry} + G^{\mathrm{wrap}}_{\bm k}
  \end{pmatrix},
  \qquad
  M_{\inn\inn} = D_{\inn\inn} + G^{\mathrm{intra}}_{\inn\inn},
  \label{eq:blockstructure}
\end{equation}
where $G^{\mathrm{intra}}$ collects the ($k$-independent) intra-cell bonds
and $G^{\mathrm{wrap}}_{\bm k}$ the wrapping bonds
\eqref{eq:wrapbond}. Crucially, $M_{\inn\inn}$, $D_{\inn\bdry}$ and
$D_{\bdry\bdry}$ are all independent of $\bm k$: intra-cell bonds never
couple to boundary modes, and $D_{\uc}$ couples modes within one site only.

Eliminating the inner modes by the Schur complement of
\eqref{eq:blockstructure} yields \emph{effective} unit-cell blocks
\begin{equation}
  \boxed{\;
  \begin{aligned}
    A_{\mathrm{eff}} &= A_{\uc} + B_\inn\, M_{\inn\inn}^{-1} B_\inn^{T} ,\\
    B_{\mathrm{eff}} &= B_\bdry - B_\inn\, M_{\inn\inn}^{-1} D_{\inn\bdry} ,\\
    D_{\mathrm{eff}} &= D_{\bdry\bdry} + D_{\inn\bdry}^{T}\, M_{\inn\inn}^{-1} D_{\inn\bdry} ,
  \end{aligned}\;}
  \label{eq:schur}
\end{equation}
such that for every momentum
\begin{equation}
  G^{(\Psi)}_{\bm k}
  = B_{\mathrm{eff}} \big( D_{\mathrm{eff}} + G^{\mathrm{wrap}}_{\bm k} \big)^{-1}
    B_{\mathrm{eff}}^{T} + A_{\mathrm{eff}} .
  \label{eq:reducedmap}
\end{equation}
Both $A_{\mathrm{eff}}$ and $D_{\mathrm{eff}}$ are real antisymmetric (using
$M_{\inn\inn}^{-T} = -M_{\inn\inn}^{-1}$), so
\eqref{eq:reducedmap} has exactly the form \eqref{eq:gaussianmap}: the
contraction of the intra-cell bonds is itself a Gaussian map, and
$(A_{\mathrm{eff}}, B_{\mathrm{eff}}, D_{\mathrm{eff}})$ are precisely the
blocks of the CM $\Guc$ of the \emph{unit-cell fiducial state} obtained by
contracting the virtual bonds inside the unit cell. 
Physically, \eqref{eq:schur} constructs
$\Guc$ from the per-site $\Gamma_s$, \eqref{eq:reducedmap} then contracts
only the bonds crossing the unit-cell boundary.

\paragraph{Complexity.} Per loss evaluation, the elimination
\eqref{eq:schur} costs one $\mathcal{O}(d_\inn^3)$ factorization
\emph{independent of the number of momenta} $N_k$, after which each
momentum requires only an $\mathcal{O}(d_\bdry^3)$ solve. The per-momentum
cost improves by a factor $\big(d / d_\bdry\big)^3 = \big(2 L_x L_y /
(L_x + L_y)\big)^3$: for a $2\times 2$ cell a factor $8$, for a
$4 \times 4$ cell a factor $64$. The memory footprint of the precomputed
bond CMs shrinks by $\big(d/d_\bdry\big)^2$ as only the wrap-bond blocks
are stored per momentum. For $L_x = L_y = 1$ every virtual mode is a
boundary mode, the inner block is empty, and
\eqref{eq:reducedmap} reduces exactly to the previous implementation, so
the trivial unit cell suffers no overhead.

\paragraph{Gradients.} The elimination \eqref{eq:schur} consists of dense
solves and products and is differentiated by automatic differentiation
(Zygote) without custom rules; it is executed once per loss evaluation. The
momentum loop \eqref{eq:reducedmap} reuses the existing hand-written
reverse rule of the Gaussian map, which stores the LU factorizations of
$D_{\mathrm{eff}} + G^{\mathrm{wrap}}_{\bm k}$ from the forward pass and is
threaded over momenta.

\end{document}
